\documentclass[10pt,a4paper]{amsart}
\usepackage[margin=19mm]{geometry}
\usepackage{amsmath,amssymb,mathtools}
\usepackage[scr=rsfs,cal=euler]{mathalfa}
\usepackage{xfrac}
\usepackage[unicode,hidelinks,bookmarks=false]{hyperref}
\newcommand{\ie}{i.e.\ }
\newcommand{\cf}{cf.\ }
\newcommand{\resp}{respectively\ }
\newcommand{\Subsection}{\subsection}
\providecommand{\nobreakdash}{\nobreak\hbox{-}}
\setlength{\emergencystretch}{2em}
\multlinegap0pt
\usepackage{amssymb}
\usepackage{breqn}
\usepackage{mathtools}
\usepackage{bm}
\usepackage[shortlabels]{enumitem}
\newtheorem{lemma}{Lemma}
\newtheorem{proposition}{Proposition}
\newcommand{\R}{\mathbb{R}}
\newcommand{\sgn}{\operatorname{sgn}}
\newcommand{\til}{\tilde}
\title{Cuspidal edges on focal surfaces of regular surfaces}
\author{Keisuke Teramoto}
\date{Source: arXiv:2509.00860v3 (CC BY 4.0)}
\begin{document}
\maketitle

\noindent\emph{Source and licence.} Cuspidal edges on focal surfaces of regular surfaces. Version arXiv:2509.00860v3 (CC BY 4.0), distributed by arXiv under the Creative Commons Attribution 4.0 International licence, http://creativecommons.org/licenses/by/4.0/, as recorded in the arXiv licence metadata for this exact version and shown on the abstract page https://arxiv.org/abs/2509.00860v3. Full licence notice: \texttt{licenses/arxiv-2509.00860-CC-BY-4.0.txt}.\par\smallskip

\noindent\textit{Editorial scope.} Counted unit: the article's proposition prop:ks-para (source0.tex lines 1117--1128). The article's complete original proof of it is delivered with it (source0.tex lines 1129--1220, one \texttt{proof} environment of 92 lines). No mathematics is written, repaired or re-derived: the statement and the proof are the article's own lines, in the article's own notation, and no retained line is substituted or re-flowed.  Retained uncounted as the source-local context the proof needs: lines 420--431; lines 446--452; lines 1037--1043.  Nothing was omitted from any retained span, so the package carries no omission record.  No retained line contains a citation, so the wrapper carries no bibliography and no bibliography entry.  No retained line contains a figure environment, an image-inclusion command or drawing code, so no image is delivered and the package declares no asset dependency.  Editorial formatting only, in addition: an AMS \texttt{amsart} preamble that restates the article's own declarations and macros used by the retained lines, namely the environments \texttt{lemma}, \texttt{proposition} on separate counters, together with the standard packages those lines need.  The retained environments print in this document as Lemma 1, Proposition 1 in order of appearance, whereas the article prints the same items with the numbering of its own full document.  The complete native source of the article as distributed in the arXiv e-print for this exact version is bundled unchanged as \texttt{sources/184197/source0.tex}.
\begin{equation}\label{eq:gauss}
	\begin{aligned}
		f_{uu}&=\Gamma^{1}_{11}f_u+\Gamma^{2}_{11}f_v+L\nu
		=\frac{E_u}{2E}f_u - \frac{E_v}{2G}f_v+L\nu,\\
		f_{uv}&=f_{vu}=\Gamma^{1}_{12}f_u+\Gamma^{2}_{12}f_v
		=\frac{E_v}{2E}f_u+\frac{G_u}{2G}f_v,\\
		f_{vv}&=\Gamma^{1}_{22}f_u+\Gamma^{2}_{22}f_v+N\nu
		=-\frac{G_u}{2E}f_u+\frac{G_v}{2G}f_v+N\nu,\\
		\nu_u&=-\kappa_1f_u,\quad 
		\nu_v=-\kappa_2f_v,
	\end{aligned}
\end{equation}

\begin{lemma}\label{lem:d-kappa}
Under the above setting, we have 
\[
\frac{E_v}{2E}=\Gamma^{1}_{12}=\frac{(\kappa_1)_{v}}{\kappa_2-\kappa_1},\quad 
\frac{G_u}{2G}=\Gamma^{2}_{12}=\frac{(\kappa_2)_{u}}{\kappa_1-\kappa_2}.
\]
\end{lemma}

\begin{equation}\label{eq:diff-c}
	\begin{aligned}
		\hat{c}'(v)&=f^t_u(c(v))u'(v)+f^t_v(c(v))
		=\left(1-\frac{\kappa_2(c(v))}{\kappa_1(c(v))}\right)f_v(c(t)),\\
		\hat{\nu}'(v)&=-\kappa_1(c(v))f_u(c(v))u'(v)-\kappa_2(c(v))f_v(c(v)),
	\end{aligned}
\end{equation}

\begin{proposition}\label{prop:ks-para}
	Let $f\colon U\to\R^3$ be a regular surface without umbilic point. 
	Assume that $f$ is parametrized by the curvature line coordinate system $(u,v)$. 
	If the parallel surface $f^t$ with $t=1/\kappa_1(p)$ \textup{(}resp. $t=1/\kappa_2(p)$\textup{)} 
	is a cuspidal edge at $p\in U$, then the singular curvature of $f^t$ at $p$ is  
	\[
	\kappa_s^t=-\frac{|\kappa_1|((\kappa_1)_v^2E-(\kappa_1)_u(\kappa_2)_uG)}
	{|(\kappa_1)_u|(\kappa_1-\kappa_2)^2G\sqrt{E}}\quad 
	\left(\text{resp.}~\kappa_s^t=-\frac{|\kappa_2|((\kappa_2)_u^2G-(\kappa_1)_v(\kappa_2)_vE)}
	{|(\kappa_2)_v|(\kappa_1-\kappa_2)^2E\sqrt{G}}\right).
	\]
\end{proposition}

\begin{proof}
	We show the case that $f^t$ with $t=1/\kappa_1(p)$ 
	is a cuspidal edge at $p$. 
	Then by the previous arguments, there exists a regular curve $c(v)=(u(v),v)$ 
	passing through $p=c(v_0)$ such that $\til{\lambda}(c(v))=0$, 
	where $\til{\lambda}(u,v)=\kappa_1(u,v)-\kappa_1(p)$. 
	We set $\hat{c}(v)=f^t(c(v))$. 
	Then by \eqref{eq:diff-c}, the second order differential $\hat{c}''(v)$ can be calculated as 
	\[
	\hat{c}''(v)=-\left(\frac{\kappa_2(c(v))}{\kappa_1(c(v))}\right)'f_v(c(v))
	+\left(1-\frac{\kappa_2(c(v))}{\kappa_1(c(v))}\right)(f_{uv}(c(v))u'(v)+f_{vv}(c(v))).
	\]
	Since $\til{\lambda}(c(v))=\kappa_1(c(v))-\kappa_1(p)=0$, we have 
	\(
	(\kappa_1)_u(c(v))u'(v)+(\kappa_1)_v(c(v))=0
	\). 
	Since $(\kappa_1)_u(p)\neq0$, it holds that 
	\begin{equation}\label{eq:du2}
		u'(v_0)=-\frac{(\kappa_1)_v(p)}{(\kappa_1)_u(p)}.
	\end{equation}
	By \eqref{eq:gauss}, Lemma \ref{lem:d-kappa} and \eqref{eq:du2}, 
	we get
	\[
	\begin{aligned}
		f_{uv}u'+f_{vv}
		&=-\frac{(\kappa_1)_v}{(\kappa_1)_u}
		\left(-\frac{(\kappa_1)_v}{\kappa_1-\kappa_2}f_u
		+\frac{(\kappa_2)_u}{\kappa_1-\kappa_2}f_v\right)
		-\frac{G_u}{2E}f_u+\frac{G_v}{2G}f_v+N\nu\\
		&=\frac{(\kappa_1)_v^2E-(\kappa_1)_u(\kappa_2)_uG}{(\kappa_1)_u(\kappa_1-\kappa_2)E}f_u
		+\left(\frac{G_v}{2G}-\frac{(\kappa_1)_v(\kappa_2)_u}{\kappa_1-\kappa_2}\right)f_v+N\nu
	\end{aligned}
	\]
	at $p$. 
	Therefore we have 
	\begin{equation}\label{eq:diff-c2}
		\hat{c}''=\frac{(\kappa_1)_v^2E-(\kappa_1)_u(\kappa_2)_uG}{\kappa_1(\kappa_1)_uE}f_u
		+Af_v+N\nu
	\end{equation}
	at $p=c(v_0)$, where $A$ is some constant. 
	By \eqref{eq:diff-c} and \eqref{eq:diff-c2}, it holds that 
	\begin{equation}\label{eq:det-para}
		\begin{aligned}
		\det(\hat{c}',\hat{c}'',\nu)
		&=\frac{(\kappa_1-\kappa_2)}{\kappa_1}\frac{(\kappa_1)_v^2E-(\kappa_1)_u(\kappa_2)_uG}{\kappa_1(\kappa_1)_uE}
		\det(f_v,f_u,\nu)\\
		&=-\frac{(\kappa_1-\kappa_2)((\kappa_1)_v^2E-(\kappa_1)_u(\kappa_2)_uG)}
		{\kappa_1^2(\kappa_1)_u}\sqrt{\frac{G}{E}}
		\end{aligned}
	\end{equation}
	at $p$. 
	On the other hand, the signed area density function $\lambda^t$ of $f^t$ is 
	\[
	\lambda^t=(1-t\kappa_1)(1-t\kappa_2)\det(f_u,f_v,\nu)
	=(1-t\kappa_1)(1-t\kappa_2)\sqrt{EG}.
	\]
	Since $\eta^t=\partial_u$ is a null vector field for $f^t$, 
	we see that 
	\[
	\eta^t\lambda^t=-t(\kappa_1)_u(1-t\kappa_2)\sqrt{EG}
	=-\frac{(\kappa_1)_u(\kappa_1-\kappa_2)}{\kappa_1^2}\sqrt{EG}
	\]
	holds at $p$. 
	Moreover, we obtain $\det(c',\eta^t)=-1$. 
	Hence we get 
	\[
	\sgn(\eta^t\lambda^t\det(c',\eta^t))
	=\sgn\left(\frac{(\kappa_1)_u(\kappa_1-\kappa_2)}{\kappa_1^2}\sqrt{EG}\right)
	=\sgn((\kappa_1)_u(\kappa_1-\kappa_2))
	\]
	at $p$. 
	Thus the singular curvature $\kappa_s^t$ is calculated as 
	\[
	\begin{aligned}
		\kappa_s^t&=-\sgn((\kappa_1)_u(\kappa_1-\kappa_2))
		\frac{\kappa_1^2|\kappa_1|}{(\kappa_1-\kappa_2)^2|\kappa_1-\kappa_2|G\sqrt{G}}
		\frac{(\kappa_1-\kappa_2)((\kappa_1)_v^2E-(\kappa_1)_u(\kappa_2)_uG)}
		{\kappa_1^2(\kappa_1)_u}\sqrt{\frac{G}{E}}\\
		&=-\sgn((\kappa_1)_u(\kappa_1-\kappa_2))
		\frac{|\kappa_1|((\kappa_1)_v^2E-(\kappa_1)_u(\kappa_2)_uG)}
		{(\kappa_1)_u(\kappa_1-\kappa_2)|\kappa_1-\kappa_2|G\sqrt{E}}\\
		&=-\frac{(\kappa_1)_u(\kappa_1-\kappa_2)}{|(\kappa_1)_u||\kappa_1-\kappa_2|}
		\frac{|\kappa_1|((\kappa_1)_v^2E-(\kappa_1)_u(\kappa_2)_uG)}
		{(\kappa_1)_u(\kappa_1-\kappa_2)|\kappa_1-\kappa_2|G\sqrt{E}}\\
		&=-\frac{|\kappa_1|((\kappa_1)_v^2E-(\kappa_1)_u(\kappa_2)_uG)}
		{|(\kappa_1)_u|(\kappa_1-\kappa_2)^2G\sqrt{E}}
	\end{aligned}
	\]
	at $p$. 
	For the case of $t=1/\kappa_2(p)$, 
	one can show in a similar way.
\end{proof}

\end{document}
